\chapter{Conceitos gerais e revisão da literatura}
\label{chp:conceitos-gerais}
Conforme explorado no capítulo \ref{chp:introducao}, um ambiente de desenvolvimento na Web é essencial para o engajamento de estudantes na área de desenvolvimento de sistemas. Com isso, nossa ferramenta Web de desenvolvimento pretende disponibilizar algoritmos e um ambiente de desenvolvimento on-line, com a possibilidade de salvar os algoritmos na núvem e permitirá validar eles. Os usuários também poderão analisar onde estão acertando e errando para que possam se auto-avaliar.

A disponibilização de ambientes de desenvolvimento na Web não é novidade. Entretanto, cada ferramenta tem sua especialidade e foco em uma área de desenvolvimento. Neste capítulo serão apresentados sistemas que seguem essa temática e que inspiram essa ferramenta.

\section{CodePen}
\label{sec:codepen}
O CodePen é um ambiente de desenvolvimento colaborativo com foco no desenvolvimento de Front-end e de código aberto. O CodePen foi desenvolvido como uma comunidade de compartilhamento de demonstrações e exemplos. Ele permite que os usuários desenvolvam sites através do navegador. Ele é ideal para expor trabalhos interativos e funcionais\cite{fiala2016collaborative}.


%Neste capítulo deve ser proporcionado o estado da arte / referencial teórico
%sobre o tema a que se refere o estudo. Um bom pesquisador não deve repetir
%trabalhos já concluídos ou que já estão em andamento. Por isso esta sessão é
%onde o autor demonstra até onde vai a pesquisa atual no campo de estudos em
%questão e estabelece as bases sobre as quais desenvolverá o estudo proposto. A
%seguir são mostrados alguns exemplos de como deve-se inserir as figuras e
%tabelas. A Figura \ref{fig:exemplo} mostra um exemplo de como inserir uma
%figura no texto. A Tabela \ref{tb:exemplo} mostra o exemplo de como uma tabela
%deve ser inserida.  Voce pode referenciar capítulos e seções adicionando labels
%à elas. Por exemplo, descrevemos a introdução no Capítulo
%\ref{chp:introducao}.

\begin{figure}[htb]
    \centering
    \includegraphics[width=0.45\textwidth]{images/figura.png}
    \caption{Exemplo de como inserir Figura}
    \label{fig:exemplo}
\end{figure}
%Altere o numero junto ao 'width' para acertar o tamanho da figura. Coloque
%numeros maior que 0 e até 1

\begin{table}[htb]
\caption{Modelo de como as tabelas devem ser inseridas no texto}
\label{tb:exemplo}
\centering
\begin{tabular}{|l|c|r|r|} %left, center, right. Você pode mudar isso
\hline
Índice  & Coluna 1 & Coluna 2 & Coluna 3 \\
\hline
Linha 1 &          &          &          \\
Linha 2 &          &          &          \\
Linha 3 &          &          &          \\
\hline
\end{tabular}
\end{table}
%Para manter a sanidade, recomenda-se deixar sempre os '&' alinhados em todas
%as linhas e colunas
